\section{Выводы для устройства памяти агента}
\label{sec:discussion}
\subsection{Содержание памяти важнее одного ограничения размера}
Результаты разделяют два требования к состоянию: ограничить активный контекст и сохранить сведения для последующих решений. Первое достигается даже бедным представлением. Второе требует выбора содержания. В поисковом сравнении записи изменений и команд соответствуют лучшим результатам, чем список файлов; добавление следующих групп сведений не даёт последовательного выигрыша. Это аргумент в пользу явного проектирования автоматической памяти, но не в пользу её неограниченного расширения.

Разделение ответственности задаёт полезный принцип такого проектирования. Захваченные результаты инструментов могут сохраняться без генерации моделью, а гипотезы и планы остаются семантическим содержанием. К этому необходимо добавить доступ к информации, которую ни один из этих способов не удержал. В испытанных конфигурациях такую роль частично выполняет недавняя история. При этом мы не установили, обеспечивает ли состояние преимущество перед тем же окном взаимодействий без него.

\subsection{Компромисс следует выбирать по качеству и затратам}
Компактное состояние не является готовой политикой экономии. При полной замене истории оно уменьшает пик контекста, но проигрывает по успешности и может увеличить стоимость сопоставимого успеха. Сохранение недавней истории меняет этот компромисс: для состояния модели оно даёт более низкую стоимость на наблюдаемый успех при меньшем точечном отставании от контроля. Для hybrid восстановление успешности оплачивается дополнительными ресурсами.

Поэтому практический выбор должен исходить из ограничения конкретной системы. Если приоритетом является успешность, потеря истории требует обоснования. Если ограничивает доступное контекстное окно, структурированная память предоставляет измеримую экономию места, но её нужно проверять на способности завершать задачи. Если важен бюджет, необходимо учитывать выход и кеш наряду с входом, а расходы оценивать вместе с неудачами. Одного отношения размеров промпта для такого решения недостаточно.

\subsection{Частичная компенсация и непроверенные продолжения}
Уже измеренная компенсация состоит в сохранении недавних взаимодействий внутри агента со состоянием. Более богатая автоматическая память представляет перспективное направление по результатам поискового сравнения; проверенные схема и inline-протокол не заменили доступ к истории. Это разные по силе выводы, и они не образуют доказательства превосходства единой новой архитектуры.

Адаптивное переключение между транскриптом и состоянием могло бы учитывать текущую потребность в контексте, а внешний архив с поиском --- возвращать наблюдения за пределами недавнего окна. Эти механизмы следуют из различия между активной памятью и сохраняемыми свидетельствами, но в работе не оценивались. Сравнение фиксированных окон не определяет момент переключения, а возврат последних взаимодействий не проверяет восстановление давно удалённых сведений. Такие вопросы требуют собственных контролируемых сравнений.
